\scanpage{175}
% L1–7
Revenant au cas d'une ext. générale
\diagram{\begin{tikzcd}[ampersand replacement=\&,column sep=small]
0\arrow[r,dash]\&E^0\arrow[r,dash]\&E\arrow[r,dash]\&\underset{\substack{\Vert\\E/E^0}}{F}\arrow[r]\&0
\end{tikzcd}}
et d'une $\lambda$-section
\[
\pi:F\longrightarrow E,
\]
associons lui
\[
\begin{aligned}
\deleted{\pi^{\unsure{t|*|\sharp|\pi}}}u\unsure{:=|:}&\,E\mathbin{\text{---}}E^0\\
\deleted{\pi^{\unsure{t|*|\sharp|\pi}}}u(x)&=\lambda x-\pi\epsilon(x)
\end{aligned}
\]
% L8–11
Elle envoie bien $E$ dans $E^0$, car
\[
\epsilon\deleted{\pi}u(x)=\lambda\underbrace{\deleted{\unsure{u}}\epsilon(x)}_y-
\epsilon\pi\underbrace{(\epsilon(x))}=\lambda y-\underbrace{\epsilon\pi(y)}_{\lambda y}=0
\]
d'autre par sa restriction à $E^0$ est $\lambda\,\mathrm{id}_{\unsure{E^0|E}}$
\[
u\mid E^0=\lambda\,\mathrm{id}_{E^0}\qquad
\begin{aligned}&\text{i.e. }u(x)=\lambda x\\&\text{si }\unsure{\alpha|x}\in E^0\end{aligned}
\]

% L12–19
C'est la version duale du $\lambda$-scindage.
On récupère $\pi$ par $u$, par la
\unsure{forme|form.|formule}
\[
\pi(\epsilon(x))=\lambda x-u(x)
\]
en notant que le deuxième \unsure{mbre|membre} se
factorise bien par $F$, car il est nul
sur $E^0$. La donnée de $\pi$ ou de $u$
est équivalente. Mais c'est $u$ qui
sera utilisé pour la théorie des $*_{i}$.
% archive sheet number: 174
